\uuid{8JmW}
\titre{Théorème de Rolle}
\auteur{bodin ; hanani ; mzari}
\organisation{exo7}
\licence{CC BY-NC-SA 4.0}
\datecreate{2020-01-20}
\niveau{L1}
\module{Analyse}
\chapitre{Dérivabilité des fonctions réelles}
\sousChapitre{Autre}
\difficulte{3}
\theme{Dérivées, théorème de Rolle}
\exo7qcmid{lille-1.derivees.026}

\contenu{
\begin{question}
\qclassification{124.00}
\qtags{motcle=théorème de Rolle}
Soit $f$ une fonction réelle continue sur $[a,b]$, dérivable sur $]a,b[$ et telle que $f(a)=f(b)=0$. Soit $\alpha \notin [a,b]$ et $\displaystyle g(x)=\frac{f(x)}{x-\alpha }$.
\begin{answers}
    \good{On peut appliquer le théorème de Rolle à $g$ sur $[a,b]$.}
    \good{Il existe $c\in ]a,b[$ tel que $\displaystyle f'(c)=\frac{f(c)}{c-\alpha}$.}
    \good{Il existe $c\in ]a,b[$ tel que la tangente à $\mathcal{C}_f$ en $c$ passe par $(\alpha ,0)$.}
    \bad{La dérivée de $g$ est $\displaystyle g'(x)=\frac{f'(x)}{(x-\alpha )^2}$.}
\end{answers}
\begin{explanations}
La fonction $g$ est continue sur $[a,b]$ et elle est dérivable sur $]a,b[$ avec
$$g'(x)=\frac{f'(x)(x-\alpha)-f(x)}{(x-\alpha)^2}.$$
De plus $g(a)=g(b)=0$. On peut donc appliquer le théorème de Rolle. Il existe alors $c\in ]a,b[$ tel que $$g'(c)=0\Leftrightarrow f'(c)(c-\alpha)-f(c)\Leftrightarrow f'(c)=\frac{f(c)}{c-\alpha}.$$
La tangente à $\mathcal{C}_f$ en $c$ passe par le point $(\alpha ,0)$.
\end{explanations}
\end{question}
}
